% =====================================================================
%  Statistical Physics - Lab 1: Microstates of a system of coins
%  Group E  -  report/main.tex
%

% =====================================================================
\documentclass[11pt,a4paper]{article}

\usepackage[utf8]{inputenc}
\usepackage[T1]{fontenc}
\usepackage{lmodern}
\usepackage[margin=2.5cm]{geometry}
\usepackage{amsmath,amssymb}
\usepackage{graphicx}
\usepackage{booktabs}
\usepackage{xcolor}
\usepackage{fancyhdr}
\usepackage[hidelinks]{hyperref}

\pagestyle{fancy}
\fancyhf{}
\lhead{Microstates of a System of Coins}
\rhead{Statistical Physics --- Lab 1}
\lfoot{UC3M}
\rfoot{Page \thepage}

% Red marker for anything that is still pending
\newcommand{\pend}[1]{\textcolor{red}{#1}}

% Include a figure if it exists, otherwise show a placeholder box
\newcommand{\safefig}[2]{%
  \IfFileExists{figures/#1}{\includegraphics[width=#2]{figures/#1}}%
  {\fbox{\parbox[c][3cm][c]{#2}{\centering\pend{missing: figures/#1}\\(run \texttt{python main.py})}}}}

% ---------------------------------------------------------------------
%  RESULTS (Part 2) - N = 100 + M, NIA digits [79, 95, 84, 54]
% ---------------------------------------------------------------------
\newcommand{\resHM}{74.6160}            % harmonic mean of the NIA digits
\newcommand{\resM}{75}                  % M = round(harmonic mean)
\newcommand{\resN}{175}                 % N = 100 + M
\newcommand{\resNhalf}{87}              % floor(N/2); k = 87 and 88 are equally probable
\newcommand{\resPa}{1.0999\times10^{-19}}   % P(k = 30)
\newcommand{\resPb}{6.0057\times10^{-2}}    % P(k = 87)
\newcommand{\resPc}{0.88682}                % P(77 <= k <= 97)
\newcommand{\resPd}{0.99809}                % P(67 <= k <= 107)
\newcommand{\resSmaxkB}{118.488}            % S_max / k_B = ln Omega(N, 87)
\newcommand{\resSmax}{1.6359\times10^{-21}} % S_max in J/K
\newcommand{\resdSekB}{-6.9091}             % ln[Omega(112)/Omega(87)]
\newcommand{\resdSe}{-9.5390\times10^{-23}}
\newcommand{\resdSfkB}{-1.0247}             % ln[Omega(97)/Omega(87)]
\newcommand{\resdSf}{-1.4148\times10^{-23}}
\newcommand{\resdSgkB}{58.657}              % ln[Omega(87)/Omega(20)]
\newcommand{\resdSg}{8.0984\times10^{-22}}
\newcommand{\restSec}{2.582\times10^{39}}   % waiting time for k = 8, s
\newcommand{\restYears}{8.18\times10^{31}}  % same, in years
\newcommand{\restAges}{5.93\times10^{21}}   % same, in ages of the Universe
\newcommand{\resRatio}{0.1576}              % P(k=87) / P(65 <= k <= 85)
\newcommand{\resNOverflow}{1028}            % largest N with float(Omega(N,N/2)) finite
\newcommand{\resLnErr}{3\times10^{-13}}     % bound on |delta ln Omega|
\newcommand{\resSErr}{4\times10^{-36}}      % k_B * resLnErr, in J/K

\title{\textbf{Microstates and Macroscopic Thermodynamics of a Binary System}\\[2pt]
\large Statistical Physics --- Laboratory Practice 1}
\author{%
  \textbf{Group E}\\
  Claudia Valiente Gutiérrez (NIA: 100524454, GitHub: @cvgggggg)\\
  Elena Valenzuela Gómez (NIA: 100522984, GitHub: @elenavga)\\
  Justin Ryan Kotrba (NIA: 100614295, GitHub: @JustinK05)\\
  Vassily Martin Schnell (NIA: 100615779, GitHub: @vasa-schnell)\\[6pt]
  \small Department of Physics, Universidad Carlos III de Madrid}
\date{\today}

\begin{document}
\maketitle
\thispagestyle{fancy}

\begin{abstract}
We study a system of $N$ independent, unbiased coins as the simplest model of
a two-state assembly. Microstates are enumerated exactly with bitwise
operations and Python generators, multiplicities $\Omega(N,k)$ are evaluated
in log-space through the log-gamma function to avoid floating-point overflow,
and the narrowing of the macrostate distribution is followed up to the
thermodynamic limit $N\to10^{23}$. For our team parameter $N=\resN$ we compute
macrostate probabilities, Boltzmann entropy changes and waiting times, and we
discuss their connection with the Second Law of Thermodynamics.
\end{abstract}

\noindent\textbf{Repository:}
\url{https://github.com/vasa-schnell/Statistical-Physics-Lab-1-Group-E}
(release \texttt{v1.0}).

% =====================================================================
\section{Introduction and objectives}
% =====================================================================
A \emph{microstate} is the full configuration of the $N$ coins (e.g.\
\texttt{HHTHT} for $N=5$), whereas a \emph{macrostate} is fixed only by the
number of heads $k$. Since every coin has two equally likely outcomes, all
$2^N$ microstates are equiprobable and the multiplicity of a macrostate is
\begin{equation}
  \Omega(N,k)=\binom{N}{k}=\frac{N!}{k!\,(N-k)!},
  \qquad \sum_{k=0}^{N}\Omega(N,k)=2^N .
  \label{eq:omega}
\end{equation}
The probability of a macrostate is $P_k=\Omega(N,k)/2^N$ and its Boltzmann
entropy is $S=k_B\ln\Omega$, with $k_B=1.380649\times10^{-23}\,\mathrm{J\,K^{-1}}$.
The objectives are: (i) enumerate microstates and find the computational
limits of brute force, (ii) obtain the macrostate distributions and their
behaviour as $N\to10^{23}$, and (iii) compute probabilities and entropy
changes for $N=100+M$ and interpret them thermodynamically.

% =====================================================================
\section{Methodology and collaborative workflow}
% =====================================================================
\subsection{Environment and reproducibility}
All code runs in a Conda environment (\texttt{statphys\_lab1}, Python 3.11,
NumPy, SciPy, Matplotlib, pytest) defined in \texttt{environment.yml}, so the
results are identical on every team member's laptop.

\subsection{Version control and repository architecture}
Development followed a feature-branch workflow on GitHub: no direct pushes to
\texttt{main}, every change was merged through a Pull Request approved by at
least one other team member. The repository is organised as
\begin{itemize}
  \item \texttt{src/}: \texttt{microstates.py} (bitwise enumeration),
        \texttt{combinatorics.py} (log-space $\Omega$ and $P_k$),
        \texttt{thermodynamics.py} (entropy, $N=100+M$, waiting times);
  \item \texttt{tests/}: \texttt{pytest} unit tests comparing the log-space
        results with exact integer combinatorics;
  \item \texttt{main.py}: runs every part and writes the figures to
        \texttt{report/figures/};
  \item \texttt{notebooks/}, \texttt{report/}: exploration and this report.
\end{itemize}
% TODO: add 1-2 sentences on how you split the work and reviewed each other's PRs.

% =====================================================================
\section{Computational analysis of microstates}
% =====================================================================
\subsection{Enumeration for $N=7$}
Each microstate is stored as an integer $s\in[0,2^N)$ whose $i$-th bit is 1
for heads, so the full space is generated with \texttt{range(1 << N)} inside a
generator (\texttt{yield}) and the number of heads is the bit count of $s$.
For $N=7$ there are $2^7=128$ microstates (Table~\ref{tab:n7}).

\begin{table}[h]
\centering
\caption{Multiplicity $\Omega(7,k)$ and a representative microstate of each macrostate.}
\label{tab:n7}
\begin{tabular}{cccr}
\toprule
Heads $k$ & Tails $7-k$ & Sample microstate & $\Omega(7,k)$\\
\midrule
7 & 0 & \texttt{HHHHHHH} & 1\\
6 & 1 & \texttt{HHHHHHT} & 7\\
5 & 2 & \texttt{HHHHHTT} & 21\\
4 & 3 & \texttt{HHHHTTT} & 35\\
3 & 4 & \texttt{HHHTTTT} & 35\\
2 & 5 & \texttt{HHTTTTT} & 21\\
1 & 6 & \texttt{HTTTTTT} & 7\\
0 & 7 & \texttt{TTTTTTT} & 1\\
\midrule
\multicolumn{3}{r}{Total} & 128\\
\bottomrule
\end{tabular}
\end{table}

\subsection{Full-enumeration wall ($2^N$)}
The cost of brute-force enumeration grows as $O(2^N)$: every extra coin
doubles the run time. Using a generator the memory stays constant
(only one state is alive at a time), whereas storing all states in a list
would need $\approx 36\,\mathrm{B}\times2^N$ of RAM. Table~\ref{tab:bench}
reports the measured times (mean $\pm$ standard deviation over 3 runs) and
peak memory. Extrapolating $t(N)=t_{\mathrm{ref}}\,2^{N-N_{\mathrm{ref}}}$, the
largest $N$ that can be fully enumerated within 5~minutes on our laptop is
$N_{\max}=24+\log_2(300/5.857)\approx29.7$, i.e.\ $N_{\max}=29$
($2^{29}\approx5.4\times10^{8}$ microstates). Each time $N$ increases by 2 the
run time is multiplied by $\approx4.0$, confirming the $2^N$ scaling.

\begin{table}[h]
\centering
\caption{Run time (mean $\pm$ standard deviation of 3 runs) and peak memory
of the full enumeration, measured on an Intel laptop (12 logical CPUs)
with Python 3.12 and \texttt{tracemalloc}.}
\label{tab:bench}
\begin{tabular}{rrccc}
\toprule
$N$ & $2^N$ & Time (s) & Peak RAM, generator & RAM if stored in a list\\
\midrule
16 & $6.55\times10^{4}$ & $0.034\pm0.007$ & 0.6~kB & $\approx 2.4$~MB\\
18 & $2.62\times10^{5}$ & $0.092\pm0.004$ & 0.6~kB & $\approx 9.4$~MB\\
20 & $1.05\times10^{6}$ & $0.366\pm0.011$ & 0.6~kB & $\approx 38$~MB\\
22 & $4.19\times10^{6}$ & $1.449\pm0.008$ & 0.6~kB & $\approx 151$~MB\\
24 & $1.68\times10^{7}$ & $5.857\pm0.037$ & 0.6~kB & $\approx 604$~MB\\
\bottomrule
\end{tabular}
\end{table}

\subsection{Macrostate distributions ($N=20,50,80,100$)}
Factorials overflow very quickly, so $\Omega$ is evaluated in log-space,
\begin{equation}
  \ln\Omega(N,k)=\ln\Gamma(N+1)-\ln\Gamma(k+1)-\ln\Gamma(N-k+1),
  \qquad \ln P_k=\ln\Omega(N,k)-N\ln2 ,
  \label{eq:lngamma}
\end{equation}
using \texttt{scipy.special.gammaln}, and only exponentiated at the end.
Figure~\ref{fig:dist} shows that the distribution is peaked at $k=N/2$; the
peak becomes higher in absolute $\Omega$ but, relative to $N$, narrower.

\begin{figure}[h]
\centering
\safefig{distributions.pdf}{0.95\textwidth}
\caption{Left: $P_k=\Omega(N,k)/2^N$ for $N=20,50,80,100$. Right: the same
curves against $k/N$; the width shrinks as $1/\sqrt{N}$.}
\label{fig:dist}
\end{figure}

\subsection{Limits of the direct binomial calculation}
Python integers have arbitrary precision, so \texttt{math.comb} itself never
overflows; the limit appears when $\Omega$ is converted to a double-precision
float (maximum $\approx1.8\times10^{308}$) to divide, plot or take logarithms.
Since $\Omega(N,N/2)\approx2^N\sqrt{2/(\pi N)}$, this happens for
$N>\resNOverflow$. Working with $\ln\Omega$ removes the problem: $\ln\Omega$
stays of order $N\ln2$, which is perfectly representable even for $N=10^{23}$.

\subsection{Thermodynamic limit ($N\to10^{23}$)}
Using Stirling's approximation, $\ln N!\approx N\ln N-N$, the binomial
distribution tends to a Gaussian of mean $\mu=N/2$ and standard deviation
$\sigma=\sqrt{N}/2$, so the relative fluctuation is
\begin{equation}
  \frac{\sigma}{\mu}=\frac{\sqrt{N}/2}{N/2}=\frac{1}{\sqrt{N}}
  \;\xrightarrow[N=10^{23}]{}\;3.2\times10^{-12}.
\end{equation}
For a macroscopic system the fraction of heads is therefore fixed at $1/2$
to about twelve significant digits (Fig.~\ref{fig:limit}), while
$S_{\max}/k_B\approx N\ln2\approx6.9\times10^{22}$.

\begin{figure}[h]
\centering
\safefig{relative_fluctuation.pdf}{0.6\textwidth}
\caption{Relative fluctuation $\sigma/\mu=1/\sqrt{N}$ as a function of the
number of coins, up to $N=10^{23}$.}
\label{fig:limit}
\end{figure}

% =====================================================================
\section{Probabilities, entropy and thermodynamics}
% =====================================================================
\subsection{Parameter assignment}
$M$ is the harmonic mean of the last two digits of the team members' NIAs,
\begin{equation}
  M=\frac{4}{\frac{1}{79}+\frac{1}{95}+\frac{1}{84}+\frac{1}{54}}=\resHM
  \;\Rightarrow\; M=\resM,\qquad N=100+M=\resN .
\end{equation}
$M$ was rounded to the nearest integer. Because $N$ is odd, $N/2=87.5$ is not
a valid number of heads; the two central macrostates $k=87$ and $k=88$ have
the same multiplicity, and we take $k_{\mathrm{mid}}=\lfloor N/2\rfloor=\resNhalf$.

\subsection{Probability calculations}
\begin{table}[h]
\centering
\caption{Macrostate probabilities for $N=\resN$.}
\label{tab:prob}
\begin{tabular}{lll}
\toprule
Case & Macrostate / range & Probability\\
\midrule
(a) & $k=30$ heads & $\resPa$\\
(b) & $k=\resNhalf$ heads (most probable) & $\resPb$\\
(c) & $77\le k\le 97$ heads ($N/2\pm10$) & $\resPc$\\
(d) & $67\le N-k\le 107$ tails ($N/2\pm20$) & $\resPd$\\
\bottomrule
\end{tabular}
\end{table}
Even the single most probable macrostate has only a 6\,\% probability, but a
window of $\pm10$ heads around it already contains 89\,\% of all microstates
and a window of $\pm20$ contains 99.8\,\%. By the head--tail symmetry,
the range in (d) is the same whether it is counted in heads or tails.
The state $k=30$, only 57 heads away from the centre, is $5\times10^{17}$
times less likely than $k=87$: as $\sigma=\sqrt{N}/2\approx6.6$, it lies about
$8.7\sigma$ below the mean (Fig.~\ref{fig:distN}).

\begin{figure}[h]
\centering
\safefig{distribution_N.pdf}{0.6\textwidth}
\caption{Macrostate distribution for $N=\resN$. Shaded: the window
$N/2\pm10$ of case (c); dashed line: $k=30$ of case (a).}
\label{fig:distN}
\end{figure}

\subsection{Entropy changes}
With $S=k_B\ln\Omega$, the entropy change of a transition $k_1\to k_2$ is
$\Delta S=S(k_2)-S(k_1)$; here we report $\Delta S_{k_1k_2}\equiv S(k_1)-S(k_2)=k_B\ln[\Omega(N,k_1)/\Omega(N,k_2)]$
as defined in the manual. The maximum entropy is
$S_{\max}=k_B\ln\Omega(\resN,\resNhalf)=\resSmax\,\mathrm{J\,K^{-1}}$
($S_{\max}/k_B=\resSmaxkB$).

\begin{table}[h]
\centering
\caption{Entropy differences $\Delta S_{k_1k_2}=S(k_1)-S(k_2)$ for $N=\resN$.}
\label{tab:entropy}
\begin{tabular}{llcc}
\toprule
Case & $k_1\to k_2$ & $\Delta S/k_B$ & $\Delta S$ (J/K)\\
\midrule
(e) & $N/2+25=112 \to 87$ (high $\to$ maximum disorder) & $\resdSekB$ & $\resdSe$\\
(f) & $N/2+10=97 \to 87$ (moderate $\to$ maximum disorder) & $\resdSfkB$ & $\resdSf$\\
(g) & $87 \to 20$ (maximum disorder $\to$ high order) & $\resdSgkB$ & $\resdSg$\\
\bottomrule
\end{tabular}
\end{table}
In (e) and (f) $\Delta S<0$ because $k_1$ is a less probable (more ordered)
macrostate than $k_{\mathrm{mid}}$; in (g) $\Delta S>0$ because the reference
state $k=20$ is far more ordered: $\Omega(87)/\Omega(20)\approx3\times10^{25}$.

\paragraph{Uncertainty.} The only source of error is floating-point rounding
in Eq.~\eqref{eq:lngamma} (machine epsilon $\epsilon\approx2.2\times10^{-16}$).
Bounding the error of each $\ln\Gamma$ term by $\epsilon$ times its value gives
$|\delta\ln\Omega|\lesssim\resLnErr$ (we checked that the actual difference
with the exact integer result is $4\times10^{-14}$), i.e.\
$\delta S\approx\resSErr\,\mathrm{J\,K^{-1}}$, completely negligible. The value
of $k_B$ is exact in the SI since 2019.

\subsection{Thermodynamic implications}
\paragraph{(h) Waiting time.} If all $N$ coins are tossed once per second, the
expected time to obtain exactly $k=8$ heads is
\begin{equation}
  t=\frac{1\,\mathrm{s}}{P_8}=\frac{2^N}{\Omega(N,8)}\,\mathrm{s}
  =\restSec\,\mathrm{s}=\restYears\ \text{years},
\end{equation}
about $\restAges$ times the age of the Universe ($1.38\times10^{10}$ years):
such an ordered state is, in practice, never observed.

\paragraph{(i) Likelihood ratio.}
\begin{equation}
  \frac{P(k=\resNhalf)}{P(M-10\le k\le M+10)}
  =\frac{\Omega(N,\resNhalf)}{\sum_{k=65}^{85}\Omega(N,k)}=\resRatio .
\end{equation}
% TODO: check with the lab manual that this is the intended comparison
% (single state at N/2 vs. the whole range M +- 10).
The ratio is smaller than one, which does not contradict the fact that
$k=\resNhalf$ is the most probable macrostate: we are comparing a
\emph{single} macrostate with a window of 21 macrostates. Since
$M+10=85$ lies only two heads below the centre, the window $[65,85]$
covers the whole left half of the peak and accumulates
$P(65\le k\le85)=0.381$, i.e.\ about 6.3 times $P_{87}=0.060$.
Per macrostate, however, the window has an average probability of
$0.381/21=0.018$, so $k=\resNhalf$ is still 3.3 times more likely than a
typical macrostate of the window, and incomparably more likely than its edge
$k=65$, which lies $3.4\sigma$ from the mean
($\Omega(87)/\Omega(65)\approx3.4\times10^{2}$).

\paragraph{(j) Second Law.} The negative values of $\Delta S$ in (e) and (f)
seem to violate the Second Law, which states that the entropy of an isolated
system never decreases. However, the Second Law is not an exact law of
mechanics but a statistical one. Every microstate is equally probable, so
the system can fluctuate into a less probable macrostate, and this is what a
negative $\Delta S$ describes. The probability of such a fluctuation is
$\Omega(k_1)/\Omega(k_{\mathrm{mid}})=e^{\Delta S/k_B}$: about
$0.36$ for the small deviation of (f), $10^{-3}$ for (e), and
$3\times10^{-26}$ for reaching $k=20$ in (g). Small fluctuations are
therefore frequent for $N=\resN$, but large ones are exponentially
suppressed, and the system spends almost all its time in the macrostates of
maximum $\Omega$. ``Entropy increases'' thus means ``the system evolves
towards the overwhelmingly most probable macrostates'', and the Second Law
becomes an effective certainty only when $N$ is macroscopic.

\paragraph{(k) Large $N$.} In the limit $N\to10^{23}$ the conclusions above
become sharp:
\begin{itemize}
  \item The probability of any \emph{single} macrostate goes to zero, even the
        most probable one: $P_{N/2}\approx\sqrt{2/(\pi N)}\approx2.5\times10^{-12}$
        for $N=10^{23}$ (cases a, b).
  \item The probability of any window of fixed \emph{relative} width around
        $k/N=1/2$ tends to one, because $\sigma/\mu=1/\sqrt{N}\to0$. Windows of
        $\pm10$ or $\pm20$ heads, however, become negligible compared with
        $\sigma\approx1.6\times10^{11}$ (cases c, d).
  \item For a deviation of fixed relative size $\delta=k/(N/2)-1$, the Gaussian
        approximation gives $\Delta S/k_B\approx-N\delta^2/2$. Its magnitude
        grows proportionally to $N$, so these fluctuations have probability
        $e^{-N\delta^2/2}\to0$ (cases e, f, g).
  \item The waiting times of (h) grow exponentially with $N$. For a
        macroscopic system an ordered state is never observed, and the
        macrostate of maximum entropy is reached with certainty. This is the
        equilibrium state of thermodynamics.
\end{itemize}

% =====================================================================
\section{Conclusions}
% =====================================================================
A system of coins captures the essential link between microscopic
configurations and macroscopic thermodynamics. Computationally, exact
enumeration of microstates is limited to $N\approx29$ within five
minutes by its $O(2^N)$ cost, although the use of bitwise operations and
generators keeps the memory constant. Direct floating-point evaluation of
$\Omega(N,k)$ fails beyond $N=\resNOverflow$, whereas log-space combinatorics
with $\ln\Gamma$ works for any $N$ with a negligible error and has been
verified against exact integer results with unit tests.

Physically, the multiplicity concentrates around $k=N/2$ with a relative
width $1/\sqrt{N}$. For our system ($N=\resN$) a window of $\pm20$ heads
around the centre already contains 99.8\,\% of the probability, and highly
ordered states require waiting times $10^{21}$ times longer than the age of
the Universe. Entropy decreases are possible but exponentially unlikely, and
in the thermodynamic limit the maximum-entropy macrostate becomes a
certainty. This is the statistical foundation of the Second Law of
Thermodynamics.

\appendix
\section{AI collaboration and audit statement}
In compliance with course policy, AI tools were used as code and writing
assistants. Every AI-generated function was audited and validated with the
\texttt{pytest} suite (log-space results against exact integer combinatorics).
The prompts, the generated code and the modifications made by the team are
logged in \texttt{AI\_STATEMENT.md} in the repository.

\end{document}
